% GENERATED_BY: oc_core_monograph_translation_draft_pipeline_v1.py
% AI_ASSISTED_TRANSLATION_DRAFT: TRUE
% THEOREM_REVIEW_STATUS: PENDING
% THEOREM_REVIEW_MODE: FORMAL_AI_GATE__CODEX_CLI_CHATGPT
% THEOREM_REVIEW_REVIEWER_ID:
% THEOREM_REVIEW_AUTHORITY_CLASS:
% THEOREM_REVIEWED_AT:
% THEOREM_REVIEW_DOSSIER_REF:
% SOURCE_LANGUAGE: EN
% TARGET_LANGUAGE: RU
% SOURCE_EN_REF: content/k_levels/k5.tex
% ================================================================
% ==== FILE: content/k_levels/k5.tex
% ================================================================

% ==============================
%  Ontology of Continua — Core
%  K-level Module: K5
%  File: content/k_levels/k5.tex
%  Status: FULLY DEFINED
% ==============================



\paragraph{Canonical tuple projection note.}
Local K-level displays in this module are projections of the canonical public tuple \(K=(\Omega,\partial\Omega,A,\Theta,P,J,C,k,M)\). When a display omits \(M\), reorders components, or uses level-indexed shorthand, the omitted embedding context is inherited from the surrounding level discussion rather than introduced as a competing definition.

\subsubsection{\texorpdfstring{$K_5$}{K_5} Обзор}
\label{sec:k5-overview}

$K_5$ — уровень, на котором \emph{электрическая возбудимость} проявляется
как новая онтологическая размерность. Граница протоклетки $K_4$ становится
\emph{электрически активной мембраной} благодаря формированию ионных
каналов с управляемой проводимостью. Это создаёт новую ось:
\[
A_{\mathrm{exc}}: \Delta V \;\mapsto\; \text{разность электрического состояния},
\]
где $\Delta V$ — трансмембранный потенциал.

Эта ось не может быть определена в $K_4$. Она представляет новый класс
несовместимых различий и, следовательно, означает реальное увеличение
размерности.

Организационно $K_5$ характеризуется:
\begin{itemize}
    \item устойчивыми ионными каналами (открытый/закрытый/дырочный/блокированный),
    \item управляемыми ионными потоками,
    \item возбудимыми циклами (возбуждение–восстановление),
    \item рефрактерной динамикой,
    \item контролем шума и пороговых ответов,
    \item самоподдерживающейся электрической организацией, тесно
          связанной с химической средой внутри.
\end{itemize}

$K_5$ — минимальный \emph{нейроподобный} континуум, даже если
настоящих нейронов пока ещё нет.

% ================================================================
\subsubsection{Пространство состояний \texorpdfstring{$\Omega(K_5)$}{\Omega(K_5)}}

Состояние $K_5$ включает:
\[
\Omega(K_5) = \{ \Delta V,\, c_i,\, g_{\mathrm{ion}},\
                P_{\mathrm{open}},\
                x(t),\
                G=(V,E),\
                I(t) \}.
\]

Элементы:
\begin{itemize}
    \item $\Delta V$ — трансмембранный потенциал,
    \item $c_i$ — ионные концентрации внутри и снаружи,
    \item $g_{\mathrm{ion}}$ — проводимости ионных каналов,
    \item $P_{\mathrm{open}}$ — вероятности открытия,
    \item $x(t)$ — внутренние переменные активации,
    \item $G=(V,E)$ — граф прото-нейронной связности для многоузловых систем
          $K_5$,
    \item $I(t)$ — внешние или внутренние токи.
\end{itemize}

Допустимость требует:
\[
|\Delta V| < \Theta_{\mathrm{volt-max}}, \quad
  c_i \in \text{допустимый диапазон}, \quad
0 \le P_{\mathrm{open}} \le 1.
\]

% ================================================================
\subsubsection{\texorpdfstring{$\partial\Omega(K_5)$}{\partial\Omega(K_5)} (граница)}

Мембрана наследует границу $K_4$ с новой структурой:
\[
\partial\Omega(K_5)
    = \partial\Omega(K_4)
       \cup \{ \text{ионные каналы с кинетикой затвора} \}.
\]

Для каждого канала $k$:
\begin{itemize}
    \item пространство состояний: открытый/закрытый/блокированный/дырочный,
    \item параметры: $g_k$, потенциал обратного течения $E_k$, переменные
          затвора,
    \item локализация на участке,
    \item взаимодействия с локальной кривизной и натяжением.
\end{itemize}

Коллапс порядка каналов или кинетики затвора уничтожает $\Omega(K_5)$.

% ================================================================
\subsubsection{Оси \texorpdfstring{$A(K_5)$}{A(K_5)}}

Минимальные оси:
\[
A(K_5) = A(K_4)\; \cup\
\{ A_{\mathrm{exc}},\; A_{\mathrm{ion}},\; A_{\mathrm{gate}} \}.
\]

Где:
\begin{itemize}
    \item $A_{\mathrm{exc}}$: ось электрического возбуждения,
    \item $A_{\mathrm{ion}}$: различия ионных видов и потоков,
    \item $A_{\mathrm{gate}}$: различия состояний затвора ионных каналов.
\end{itemize}

Эти оси несовместимы с чисто химическими различиями $K_4$ и
означают континуум более высокой размерности.

% ================================================================
\subsubsection{Потенциалы \texorpdfstring{$P(K_5)$}{P(K_5)}}

Потенциалы включают:
\[
P(K_5)=\{
E_{\mathrm{ion}},\;
P_{\mathrm{gate}},\;
P_{\mathrm{exc}},\;
P_{\mathrm{noise}},\;
P_{\mathrm{refrac}},\;
P_{\mathrm{coupling}}
\}.
\]

Интерпретация:
\begin{itemize}
    \item $E_{\mathrm{ion}}$ — потенциалы Нернста для каждого ионного вида,
    \item $P_{\mathrm{gate}}$ — потенциалы активации затвора,
    \item $P_{\mathrm{exc}}$ — пороги возбуждения для генерации спайков,
    \item $P_{\mathrm{noise}}$ — потенциалы подавления шума,
    \item $P_{\mathrm{refrac}}$ — потенциалы восстановления рефрактерного
          состояния,
    \item $P_{\mathrm{coupling}}$ — потенциалы электрического связанного
          поведения множественных единиц через $G$.
\end{itemize}

% ================================================================
\subsubsection{Пороги \texorpdfstring{$\Theta(K_5)$}{\Theta(K_5)}}

Структура порогов:
\[
\Theta(K_5)=\{
\Theta_{\mathrm{exc}},\;
\Theta_{\mathrm{noise}},\;
\Theta_{\mathrm{volt-max}},\;
\Theta_{\mathrm{channel-open}},\;
\Theta_{\mathrm{channel-close}},\;
\Theta_{\mathrm{channel-noise}},\;
\Theta_{\mathrm{front}},\;
\Theta_{\mathrm{refrac}}
\}.
\]

Значение:
\begin{itemize}
    \item $\Theta_{\mathrm{exc}}$: порог возбуждения для генерации спайков,
    \item $\Theta_{\mathrm{noise}}$: шумовой предел, после которого возбуждение
          теряется,
    \item $\Theta_{\mathrm{volt-max}}$: максимальное безопасное
          трансмембранное напряжение,
    \item $\Theta_{\mathrm{channel-open}}$, $\Theta_{\mathrm{channel-close}}$:
          пороги жизнеспособности затвора,
    \item $\Theta_{\mathrm{channel-noise}}$: устойчивость шума в затворе,
    \item $\Theta_{\mathrm{front}}$: минимальные условия для фронтов
          распространения в многоузловых системах,
    \item $\Theta_{\mathrm{refrac}}$: минимальный уровень рефрактерного
          восстановления.
\end{itemize}

Превышение $\Theta_{\mathrm{volt-max}}$ или
$\Theta_{\mathrm{channel-noise}}$ аннигилирует континуум.

% ================================================================
\subsubsection{Потоки \texorpdfstring{$J(K_5)$}{J(K_5)}}

Потоки:
\begin{itemize}
    \item ионные токи:
    \[
    J_{\mathrm{ion},k} = g_k(\Delta V - E_k),
    \]
    \item токи утечки и шумовые токи,
    \item потоки переменных затвора:
          $\dot{x} = f(x,\Delta V)$,
    \item потоки распространения спайков по рёбрам графа $E$,
    \item потоки восстановления рефрактерности,
    \item потоки связи между единицами.
\end{itemize}

Стабильность требует:
\[
\sigma(J) < \sigma_{\max}(K_5).
\]

% ================================================================
\subsubsection{Циклы \texorpdfstring{$C(K_5)$}{C(K_5)}}

Характерные циклы:
\[
C(K_5)
= \{ C_{\mathrm{exc}},\;
     C_{\mathrm{refrac}},\;
     C_{\mathrm{noise-control}},\;
     C_{\mathrm{channel-dynamics}},\;
     C_{\mathrm{coupling}} \}.
\]

Описания:
\begin{itemize}
    \item $C_{\mathrm{exc}}$: цикл возбуждения (покой → возбуждение →
          пик → реполяризация),
    \item $C_{\mathrm{refrac}}$: рефрактерный цикл восстановления состояний
          канала,
    \item $C_{\mathrm{noise-control}}$: подавление стохастических
          флуктуаций,
    \item $C_{\mathrm{channel-dynamics}}$: циклы затвора, поддерживающие
          возбудимость,
    \item $C_{\mathrm{coupling}}$: многопоточнные синхронизированные
          циклы.
\end{itemize}

Каждый цикл имеет характерный период $\tau_{C_j}$.

% ================================================================
\subsubsection{Время \texorpdfstring{$\tau(K_5)$}{\tau(K_5)}}

Время возникает из существования устойчивых циклов возбуждения:
\[
\tau(K_5)
    = \min_j \tau_{C_j},
\qquad
\Pi(C_j) > \Theta_{\mathrm{time}}.
\]

Временной порядок в $K_5$ жёстче, чем в $K_4$, из-за резких порогов
возбуждения и рефрактерных интервалов.

% ================================================================
\subsubsection{Континуальность \texorpdfstring{$k(K_5)$}{k(K_5)}}

По общей формуле:
\[
k_5
= H(\Omega(K_5))
  \cdot\frac{|V_{\mathrm{cycle}}|}{|V|}
  \cdot\frac{\sum_j C_j^{\mathrm{eff}}}{\sum_j C_j^{\max}}
  \cdot\left(1-\frac{\sigma(J)}{\sigma_{\max}}\right)
  \cdot\left(1-\frac{T(K_5)}{\Theta_{\mathrm{stab}}}\right)_+.
\]

Интерпретация:
\begin{itemize}
    \item доля цикла управляет жизнеспособностью,
    \item избыточный шум снижает $k_5$,
    \item слишком большое $\Delta V$ обнуляет $k_5$,
    \item координированное возбуждение повышает $k_5$.
\end{itemize}

% ================================================================
\subsubsection{Структурное напряжение \texorpdfstring{$T(K_5)$}{T(K_5)}}

Источники:
\begin{itemize}
    \item несовместимые ионные градиенты,
    \item шум затвора,
    \item разгоняющее возбуждение,
    \item недостаточное восстановление рефрактерности,
    \item потеря стабильности связности,
    \item неконтролируемая деполяризация.
\end{itemize}

Сбой при:
\[
T(K_5) > \Theta_{\mathrm{volt-max}}.
\]

% ================================================================
\subsubsection{Энергия \texorpdfstring{$E(K_5)$}{E(K_5)}}

Энергетический функционал:
\[
E(K_5)
= E_{\mathrm{chem}}
+ E_{\mathrm{ion}}
+ E_{\mathrm{pump}}
+ E_{\mathrm{exc}}
+ E_{\mathrm{noise}}
+ E_{\mathrm{coupling}}.
\]

Компоненты:
\begin{itemize}
    \item химическая энергия для поддержания градиентов,
    \item энергия ионных токов,
    \item энергетические затраты насосов,
    \item энергия возбуждения и затвора,
    \item энергия подавления шума,
    \item энергия связности в многоузловых сетях.
\end{itemize}

% ================================================================
\subsubsection{Операторы на \texorpdfstring{$K_5$}{K_5} (\texorpdfstring{$\Psi$}{\Psi}, \texorpdfstring{$\Phi$}{\Phi}, \texorpdfstring{$\Lambda$}{\Lambda}, \texorpdfstring{$U$}{U}, \texorpdfstring{$\Chi$}{\Chi})}

\begin{itemize}
    \item $\Psi_{5\to6}$ — возникновение когнитивных моделирующих осей и
          превращение циклов возбуждения в аттракторы,
    \item $\Phi$ — эволюция кинетики каналов и проводимости,
    \item $\Lambda$ — структурная связь возбуждения с химической
          гомеостатикой,
    \item $U$ — унификация ионных видов в интегрированную электрическую
          динамику,
    \item $\Chi$ — ветвление в различные режимы возбудимости
          (burst, tonic, phasic).
\end{itemize}

% ================================================================
\subsubsection{Процессы на \texorpdfstring{$K_5$}{K_5}}

Типичные процессы:
\begin{itemize}
    \item запуск и восстановление спайков,
    \item переходы кинетики затвора,
    \item флуктуации вокруг порога возбуждения,
    \item распространение фронтов по $G$,
    \item рефрактерно-индуцированный порядок активности,
    \item синхронизация в связанных единицах,
    \item возникновение устойчивых паттернов фейрингов, похожих на
          аттракторы.
\end{itemize}

% ================================================================
\subsubsection{Прогнозы для \texorpdfstring{$K_5$}{K_5}}

\begin{enumerate}
    \item Возбудимость требует строго контролируемых режимов шума.
    \item Спайкоподобные события возникают спонтанно после пересечения
          $\Theta_{\mathrm{exc}}$.
    \item Глобальная деполяризация вызывает структурный коллапс.
    \item Связанные системы $K_5$ формируют аттракторы, предвосхищающие
          когнитивные модели $K_6$.
    \item Разнообразие каналов расширяет допустимые области $\Omega(K_5)$.
\end{enumerate}

% ================================================================
\subsubsection{Эксперименты для \texorpdfstring{$K_5$}{K_5}}

Эмпирические аналоги:
\begin{itemize}
    \item искусственные липидные везикулы с реконституированными ионными
          каналами,
    \item patch-clamp эксперименты измерения шума затворов,
    \item вольтамперометрические измерения порогов $\Delta V$,
    \item мульти-везикулярная связь и тесты синхронизации,
    \item протонейронные модели с минимальной возбудимостью.
\end{itemize}

Эти тесты напрямую исследуют $\Theta_{\mathrm{exc}}$, $\Theta_{\mathrm{noise}}$,
$\Theta_{\mathrm{channel-open}}$ и т.д.

% ================================================================
\subsubsection{Коллапс и смерть \texorpdfstring{$K_5$}{K_5}}

Происходит, когда:
\[
\Omega(K_5) = \varnothing.
\]

Механизмы:
\begin{itemize}
    \item неконтролируемая деполяризация,
    \item неисправность канала или коллапс кинетики затвора,
    \item рассасывание ионных градиентов,
    \item шум-индуцированные неконтролируемые отказы,
    \item неспособность восстанавливаться после возбуждения.
\end{itemize}

Коллапс препятствует возникновению $K_6$.

% ================================================================
\subsubsection{Фальсифицируемость \texorpdfstring{$K_5$}{K_5}}

Основные проверяемые прогнозы:
\begin{enumerate}
    \item Возбудимое поведение не может сохраняться без рефрактерных
          циклов.
    \item Шумовые пороги управления реально существуют и измеримы.
    \item Распределения ионных каналов определяют границы жизнеспособности.
    \item Устойчивые аттракторы в $K_5$ должны существовать до формирования
          $K_6$.
    \item Для распространения спайков требуется $\Theta_{\mathrm{front}}$.
\end{enumerate}

Отрицание этих прогнозов опровергает модель $K_5$.

% ================================================================
\subsubsection{Разветвление / Онтологическое положение \texorpdfstring{$K_5$}{K_5}}

Типы ветвления:
\begin{itemize}
    \item режимы возбудимости tonic–phasic–burst,
    \item многоканальные и одно-канальные системы,
    \item слабо и сильно связанные единицы,
    \item распределённые и локализованные паттерны возбуждения.
\end{itemize}

Онтологическое положение:
\[
K_4 \to K_5 \to K_6.
\]

$K_5$ — мост между биологическими протоклетками и когнитивными системами.

% ================================================================
\subsubsection{Связь с M-пространствами}

$M_5$ задаёт:
\begin{itemize}
    \item разрешённые ионные виды и носители заряда,
    \item допустимые диапазоны проводимости и потенциалов,
    \item шумовые режимы, совместимые с возбудимостью,
    \item ограничения среды, позволяющие устойчивость каналов,
    \item разрешённые структуры связности.
\end{itemize}

Совместимость $K_5$ с $M_5$ определяет существование
электрически возбудимых систем.
